% ============================================================
% clock_fluid_dynamics.tex
% STATUS: STUB -- theoretical development in progress
% See notes/clock_fluid_dynamics.md for working notes
% ============================================================

% This section derives GR as the continuum limit of clock fluid dynamics.
% It is the central theoretical result of v2.0 if the derivation succeeds.

% Topics to cover:

% 7.1 The Global Clock Budget
% ---------------------------------------------------------
% A=1 applied universally implies a conserved total clock count.
% The clock density field rho_clock(x,y,z,t) is the fundamental
% gravitational field -- not a metric, not a potential, but a density.

% 7.2 The Continuity Equation
% ---------------------------------------------------------
% d(rho_clock)/dt + div(J_clock) = 0
% Derived from CausalSession.tick() in the continuum limit.
% Physical interpretation: clocks are conserved, not created or destroyed.
% Consequence: gravitational time dilation has the correct sign from
% clock conservation alone -- no geometry required.

% 7.3 The Momentum Equation
% ---------------------------------------------------------
% d(J)/dt + (J.grad)J = -(c^2/rho)*grad(rho) + (V(rho)-J)*omega
% Derived from the phase-alignment weighting in tick().
% The pressure term -(c^2/rho)*grad(rho) is the refractive bias.
% The relaxation term (V(rho)-J)*omega encodes inertia correctly:
%   high omega -> fast relaxation -> low inertia resistance (to investigate)
%   low omega  -> slow relaxation -> high inertia resistance
% Connection to Payne-Whitham traffic flow model.

% 7.4 The Planck Density Ceiling
% ---------------------------------------------------------
% rho_max = 1 / l_P^3: maximum clock density at Planck floor.
% Physical consequences:
%   - No singularities: density bounded, not infinite
%   - Event horizon = scheduler saturation (rho -> rho_max)
%   - Natural UV cutoff replacing GR renormalization problems
% Bekenstein-Hawking entropy as boundary clock count at saturation:
%   S = A / (4 l_P^2)  [factor of 4 from lattice geometry -- TBD]

% 7.5 The Superfluid Limit
% ---------------------------------------------------------
% Under A=1 the clock fluid is:
%   - Conservative (zero viscosity -- no amplitude lost)
%   - Compressible (rho_clock varies)
%   - Quantized (U(1) phase gives quantized vorticity)
%   - Density-bounded (Planck floor)
% This is the equation of state of a compressible superfluid.
% Gravitational waves as phonons in the clock fluid:
%   omega_gw = c_eff * |k| at low density (matches GR)
%   Nonlinear dispersion near rho_max: testable prediction.

% 7.6 GR as the Continuum Limit (the main result)
% ---------------------------------------------------------
% In the limit N -> infinity (node spacing -> 0, tick -> 0):
%   the clock fluid equations reduce to the Einstein field equations
%   with g_mu_nu = f(rho_clock, J_clock).
% The discrete corrections at finite N are the falsifiable predictions
% (connected to exp_06 path counting results).

% 7.7 Black Hole Formation as a Clock Fluid Shock Wave
% ---------------------------------------------------------
% The traffic flow analogy predicts shock waves (discontinuities in rho).
% In the clock fluid: a shock wave in rho_clock is a black hole forming.
% The shock front propagates at c (the maximum signal speed).
% Inside the shock: rho -> rho_max, time halts.
% Outside: rho drops sharply, time runs fast (gravitational redshift).

\textbf{[STUB -- see notes/clock\_fluid\_dynamics.md]}

% Key references for this section:
% Traffic flow analogy: \cite{lighthill1955} \cite{payne1971}
% Bekenstein-Hawking entropy: \cite{bekenstein1973} \cite{hawking1975}
% Holographic principle: \cite{maldacena1998}
% Entropic gravity comparison: \cite{verlinde2011}
% Information entropy: \cite{shannon1948}
